% Related Work Section
% Auto-generated by conceptgraph.evaluation.related_work
% Generated: 2026-03-20 05:48:45
%
% This file compares our two-stage evidence-seeking framework with related methods.
% Key differentiation axes:
% 1. Adaptive Evidence Acquisition
% 2. Symbolic-to-Visual Repair
% 3. Evidence-Grounded Uncertainty
% 4. Unified Multi-Task Policy

\section{Related Work}
\label{sec:related}

We review prior work on 3D scene understanding with vision-language models, scene graphs for embodied AI, iterative reasoning approaches, and multi-task embodied agents. Our method distinguishes itself through \emph{adaptive evidence acquisition}, \emph{symbolic-to-visual repair}, \emph{evidence-grounded uncertainty}, and \emph{unified multi-task policy}.

\subsection{3D Scene Understanding with Vision-Language Models}

Recent advances in vision-language models have enabled impressive 3D scene understanding capabilities~\cite{3dgraphllm2025,scenevlm2025,openground2025}. 3DGraphLLM~\cite{3dgraphllm2025} represents scene graphs as learnable token sequences that LLMs consume directly, achieving strong performance on ScanRefer (62.4\% Acc@0.25). However, this approach treats scene graphs as fixed input representations---errors in detection propagate unchanged to the output.

Scene-VLM~\cite{scenevlm2025} introduces iterative feedback loops with specialized modules for generation, transformation, and judgment. While effective for layout optimization, the iteration pattern is predetermined rather than adaptive to task requirements. OpenGround~\cite{openground2025} uses active cognition-based reasoning to overcome limitations of pre-defined object tokens, but focuses narrowly on visual grounding without broader task coverage.

These methods demonstrate the value of combining 3D structure with language models, but lack mechanisms for recovering from detection failures or dynamically acquiring additional evidence when initial inputs are insufficient. Our framework addresses these gaps through Stage 2's adaptive evidence-seeking policy.

\subsection{Scene Graphs for Embodied AI}

Scene graphs provide structured representations of spatial relationships that benefit embodied reasoning~\cite{sgnav2024,ovigo2025}. SG-Nav~\cite{sgnav2024} constructs online scene graphs during navigation, using hierarchical chain-of-thought reasoning for zero-shot object navigation. OVIGo-3DHSG~\cite{ovigo2025} extends this with multi-level hierarchical graphs (floors, rooms, locations, objects) and LLM-guided multi-hop reasoning.

A key limitation of these approaches is that scene graphs are treated as ground truth. When object detection fails or relationships are misclassified, downstream reasoning has no mechanism for correction. In contrast, our framework treats Stage 1's scene graph hypotheses as \emph{soft priors}---Stage 2 can validate, refine, or reject these hypotheses through visual evidence, enabling recovery from detection failures that would be catastrophic for pure graph-based methods.

\subsection{Iterative Reasoning in Vision-Language Models}

Several recent works explore multi-step reasoning for visual understanding~\cite{probeground2026,scenevlm2025}. Most relevant is the concurrent work on evidence-seeking agents~\cite{probeground2026}, which trains an RL policy for Probe-and-Ground actions that dynamically acquire visual evidence. This approach validates our core claim that iterative evidence-seeking outperforms non-interactive baselines (they report 18\% Brier score improvement).

Our framework differs in three key aspects: (1) we operate on 3D scene understanding with explicit scene graph structure, while~\cite{probeground2026} focuses on 2D diagnostic reasoning; (2) we provide symbolic-to-visual repair that corrects scene graph errors, not just evidence accumulation; (3) we support a unified policy across QA, grounding, navigation, and manipulation rather than single-task optimization.

The ReAct paradigm~\cite{react2022} provides a general framework for combining reasoning and acting that we build upon. Our contribution is adapting this paradigm specifically for embodied 3D scene understanding with task-conditioned evidence refinement and structured hypothesis repair.

\subsection{Multi-Task Embodied Agents}

Unified agents that handle multiple embodied tasks have gained significant attention~\cite{leo2024}. LEO~\cite{leo2024} demonstrates a generalist agent capable of 3D captioning, question answering, navigation, and manipulation using dual encoders and two-stage training. While LEO achieves broad task coverage, it requires extensive task-specific training and processes inputs in a single forward pass without dynamic evidence refinement.

Recent work explores unified token representations~\cite{rt2} and foundation models~\cite{embodnav} for multi-task embodied AI. These approaches focus on shared representations and joint training rather than task-conditioned evidence seeking.

Our unified policy takes a complementary approach: rather than training a single model end-to-end, we provide Stage 2 with tools that adapt evidence acquisition to task requirements. This design enables strong performance across tasks without task-specific training, as the same evidence-seeking strategy (requesting views, crops, or hypothesis changes) benefits QA, grounding, navigation, and manipulation.

Table~\ref{tab:related-work-comparison} summarizes the comparison with related methods across key dimensions.

\begin{table*}[t]
\centering
\small
\renewcommand{\arraystretch}{1.15}
\setlength{\tabcolsep}{4pt}
\begin{tabular}{l c c c c c c c}
\toprule
Method & Venue & \makecell{Evidence\\Acquisition} & \makecell{Scene\\Graph} & \makecell{Detection\\Recovery} & \makecell{Uncertainty\\Output} & \makecell{Multi-Task\\Unified} & Tasks \\
\midrule
3DGraphLLM & ICCV 2025 & Static & \cmark & \xmark & \xmark & \cmark & C/G/Q \\
SG-Nav & NeurIPS 2024 & Iterative & \cmark & \xmark & \xmark & \xmark & N \\
Scene-VLM & arXiv 2025 & Iterative & \xmark & \xmark & \xmark & \xmark & G/Q \\
LEO & ICML 2024 & Static & \xmark & \xmark & \xmark & \cmark & C/M/N/Q \\
OpenGround & CVPR 2025 & Iterative & \xmark & \xmark & \xmark & \xmark & G \\
Probe-and-Ground & CVPR 2026 & \textbf{Adaptive} & \xmark & \xmark & \cmark & \xmark & Q \\
OVIGo-3DHSG & arXiv 2025 & Iterative & \cmark & \xmark & \xmark & \xmark & N \\
\midrule
\textbf{Two-Stage Evidence-Seeking Agent} & arXiv 2026 & \textbf{Adaptive} & \cmark & \cmark & \cmark & \cmark & G/M/N/Q \\
\bottomrule
\end{tabular}
\caption{Comparison of our two-stage evidence-seeking framework with related methods. Tasks: Q=QA, G=Grounding, N=Navigation, M=Manipulation, C=Captioning. Our method uniquely combines adaptive evidence acquisition, detection failure recovery, uncertainty quantification, and unified multi-task support.}
\label{tab:related-work-comparison}
\end{table*}